\section{Machine Learning}
\label{sec:ml}

The project includes historical Python prototypes and Rust signal code. The
legacy ML archive includes XGBoost, LightGBM, logistic meta-classifiers, SHAP
feature analysis, and stacked ensembles. The \textbf{published} binary-outcome
model uses a Rust pipeline on unified Parquet:

\begin{enumerate}
  \item \texttt{export\_step3\_from\_parquet} reads \texttt{v0.4.2-unified}
        tables and emits step3 calibration CSVs (43 features).
  \item \texttt{train\_binary\_outcome\_model} performs walk-forward logistic
        regression by market, applies Platt scaling, and writes JSON artifacts.
\end{enumerate}

On the publication workstation (Apple M5 Max), step3 export of 354{,}684 rows
across 2{,}234 markets completes in ${\sim}$78\,s; training (555 walk-forward
windows) completes in ${\sim}$67\,s. Step3 export writes 50\% of
\texttt{market\_meta} markets; the remainder lack sufficient Polymarket ticks
or Binance trades.

For public releases, model binaries are uploaded to Hugging Face Models rather
than committed to Git. The recommended \texttt{v0.2/} release achieves
calibrated walk-forward AUC 0.840, Brier 0.164, and ECE 0.025 on 354k step3
rows. Simulated positive-EV trading under 1\% fees and 0.5\% slippage yields
$-0.123$ PnL per trade---reported as a reproducibility baseline, not deployable
alpha. An earlier \texttt{v0.1/} pilot artifact remains for comparison.

\begin{figure}[t]
  \centering
  \begin{tikzpicture}[
      node distance=0.4cm,
      every node/.style={flowbox, minimum width=2.9cm}
    ]
    \node (pq) {Unified Parquet};
    \node[below=of pq] (s3) {Step3 export};
    \node[below=of s3] (split) {Walk-forward split};
    \node[below=of split] (fit) {Logistic + Platt};
    \node[below=of fit] (exp) {HF artifact};
    \draw[flowarr] (pq) -- (s3);
    \draw[flowarr] (s3) -- (split);
    \draw[flowarr] (split) -- (fit);
    \draw[flowarr] (fit) -- (exp);
  \end{tikzpicture}
  \caption{Published ML pipeline: Parquet-native step3 export, walk-forward
    training, external artifact storage.}
  \label{fig:training}
\end{figure}